\chapter{Naive Bayesian classification}
\label{chap:naive-bayes}

Naive Bayes is a probabilistic classifier based on Bayes' theorem. It is called
\emph{Naive} because it assumes that features are conditionally independent given
the class label. Despite this strong assumption, it performs well in various
applications, including text classification, spam filtering, and sentiment
analysis.

\begin{warning}
In practice the hypothesis (see~\eqref{eq:ch6-naive-assumption}) of the naive
Bayes model is seldom exactly satisfied. It can nevertheless be used as a
baseline and has surprisingly good results; this means that a model that
performs worse than naive Bayes should probably not be used.
\end{warning}

\section{Worked-out example: Bag-of-Words sentiment analysis}
\label{sec:ch6-bow-example}

We will start with a worked-out example to apply the Naive Bayes classifier
to a simple sentiment analysis problem. Given a collection of labeled reviews
composed of words, we will compute probabilities to classify a new review as
`positive' ($+$) or `negative' ($-$).

Bayes' theorem states:
\begin{equation}\label{eq:ch6-bayes}
  \PP(C\mid X)=\frac{\PP(X\mid C)\,\PP(C)}{\PP(X)}
\end{equation}
where:
\begin{itemize}
  \item $\PP(C\mid X)$ is the posterior probability of class
        $C\in\{+,-\}$ given the data (the review) $X$.
  \item $\PP(X\mid C)$ is the likelihood of data $X$ given class $C$.
  \item $\PP(C)$ is the prior probability of class $C$.
  \item $\PP(X)$ is the evidence (which remains the same for all classes)
        assumed to be strictly positive.
\end{itemize}

Here the `Review' is composed of words denoted $w$. The $\PP(X)$ is usually
unknown but the same for all classes so we only need to compare:
\begin{align}
  \PP(\text{Positive}\mid\text{Review})
    &\propto \PP(\text{Review}\mid\text{Positive})\,\PP(\text{Positive}),
      \label{eq:ch6-prop-pos}\\
  \PP(\text{Negative}\mid\text{Review})
    &\propto \PP(\text{Review}\mid\text{Negative})\,\PP(\text{Negative}).
      \label{eq:ch6-prop-neg}
\end{align}

The \emph{Naive Bayes assumption} states that words contribute independently
to the probability, leading to:
\begin{equation}\label{eq:ch6-bow-product}
  \PP(\text{Review}\mid C)=\prod_{w\in\text{Review}}\PP(w\mid C).
\end{equation}

\begin{remark}
The procedure here belongs to the general framework of `Bag-of-Words' (BoW)
approaches which treat the text as an unordered collection of words,
disregarding grammar and word order while preserving word frequency.
\end{remark}

We have a dataset of labeled movie reviews (Table~\ref{tab:ch6-training}).

\begin{table}[ht]
  \centering
  \begin{tabular}{ll}
    \toprule
    \textbf{Review} & \textbf{Sentiment}\\
    \midrule
    love this movie    & Positive ($+$)\\
    hate this film     & Negative ($-$)\\
    great acting love  & Positive ($+$)\\
    boring and slow    & Negative ($-$)\\
    amazing film       & Positive ($+$)\\
    bad and boring     & Negative ($-$)\\
    \bottomrule
  \end{tabular}
  \caption{Training Data}
  \label{tab:ch6-training}
\end{table}

We want to classify the review:
\begin{equation}\label{eq:ch6-testreview}
  \text{``love boring movie''}.
\end{equation}

\paragraph{Step 1: Compute Priors}
\begin{equation}\label{eq:ch6-priors}
  \PP(+)=\frac{3}{6}=0.5,\qquad \PP(-)=\frac{3}{6}=0.5.
\end{equation}

\paragraph{Step 2: Compute Likelihoods with Laplace Smoothing}
The collection of all words, called ``vocabulary'', is:
\begin{align}
  V&=\{\text{love, this, movie, hate, film, great, acting, boring,
        and, slow, amazing, bad}\},\label{eq:ch6-vocab}\\
  |V|&=12.\label{eq:ch6-vocabsize}
\end{align}

The word frequencies are given in Table~\ref{tab:ch6-wordcounts}.

\begin{table}[ht]
  \centering
  \small
  \begin{tabular}{lcccccccccccc}
    \toprule
      & \rotatebox{60}{love} & \rotatebox{60}{this} & \rotatebox{60}{movie}
      & \rotatebox{60}{hate} & \rotatebox{60}{film} & \rotatebox{60}{great}
      & \rotatebox{60}{acting} & \rotatebox{60}{boring} & \rotatebox{60}{and}
      & \rotatebox{60}{slow} & \rotatebox{60}{amazing} & \rotatebox{60}{bad}\\
    \midrule
    Positive ($+$) & 2 & 1 & 1 & 0 & 1 & 1 & 1 & 0 & 0 & 0 & 1 & 0\\
    Negative ($-$) & 0 & 1 & 0 & 1 & 1 & 0 & 0 & 2 & 2 & 1 & 0 & 1\\
    \bottomrule
  \end{tabular}
  \caption{Word Counts per Class}
  \label{tab:ch6-wordcounts}
\end{table}

We now compute conditional probabilities. Note that in our procedure
these probabilities will be multiplied; an obvious problem appears if one is
zero because it will make the whole product null\footnote{This is a problem
because, for instance, it would mean any word $w_1$ that never appeared
in a Positive ($+$) labeled input has $\PP[w_1\mid(+)]=0$; also any word
$w_2$ that never appeared in a Negative ($-$) labeled input has
$\PP[w_2\mid(-)]=0$. But for a text $t$ containing both $w_1$ and $w_2$
we would have $\PP[t\mid(+)]=\PP[t\mid(+)]=0$!}. To address this problem we
use Laplace smoothing. Laplace smoothing is a technique used in probabilistic
models to handle zero probabilities by adding a small constant $\alpha$
(typically $1$) to all counts, modifying the probability estimate as
\begin{equation}\label{eq:ch6-laplace-general}
  \PP(w_i\mid C)=\frac{\mathrm{Count}(w_i,C)+\alpha}
                      {\sum_{w}\mathrm{Count}(w,C)+\alpha|V|},
\end{equation}
ensuring that no probability is ever exactly zero. Using Laplace smoothing
($\alpha=1$):
\begin{equation}\label{eq:ch6-laplace-alpha1}
  \PP(w_i\mid C)=\frac{\mathrm{Count}(w_i,C)+1}
                      {\sum_{w}\mathrm{Count}(w,C)+|V|}.
\end{equation}

For the positive class ($+$):
\begin{align}
  \PP(\text{love}\mid +) &= \frac{2+1}{8+12}=\frac{3}{20}=0.15,
  & \PP(\text{boring}\mid +) &= \frac{0+1}{8+12}=\frac{1}{20}=0.05,
    \label{eq:ch6-likel-pos-1}\\
  \PP(\text{movie}\mid +) &= \frac{1+1}{8+12}=\frac{2}{20}=0.10. & &
    \label{eq:ch6-likel-pos-2}
\end{align}

For the negative class ($-$):
\begin{align}
  \PP(\text{love}\mid -) &= \frac{0+1}{9+12}=\frac{1}{21}\approx 0.048,
  & \PP(\text{boring}\mid -) &= \frac{2+1}{9+12}=\frac{3}{21}=0.143,
    \label{eq:ch6-likel-neg-1}\\
  \PP(\text{movie}\mid -) &= \frac{0+1}{9+12}=\frac{1}{21}\approx 0.048. & &
    \label{eq:ch6-likel-neg-2}
\end{align}

\paragraph{Step 3: Compute Posterior Probabilities}
For Positive ($+$):
\[
  \PP(\text{Positive}\mid\text{Review})\propto
  0.5\times 0.15\times 0.05\times 0.10 = 0.000375.
\]
For Negative ($-$):
\[
  \PP(\text{Negative}\mid\text{Review})\propto
  0.5\times 0.048\times 0.143\times 0.048 = 0.000164.
\]

\paragraph{Step 4: Compare and Classify}
Since
\[
  \PP(\text{Positive}\mid\text{Review})=0.000375 >
  \PP(\text{Negative}\mid\text{Review})=0.000164,
\]
we classify review ``love boring movie'' as Positive ($+$). Despite its
simplifying assumptions, Naive Bayes is highly effective for text
classification tasks like spam filtering and sentiment analysis. Using
probability estimates and Laplace smoothing, we were able to determine that
``love boring movie'' should be classified as a positive review.

\section{Statistical Formulation of Naive Bayes}
\label{sec:ch6-statistical}

\subsection{\emph{Discriminative Model} vs \emph{Generative Model}}
\label{subsec:ch6-disc-vs-gen}

\begin{definition}\label{def:ch6-disc-gen}
\begin{enumerate}
  \item A \emph{discriminative model} is a model for the conditional
        probability $p(y\mid x):=\PP(Y=y\mid X=x)$ where
        $(x,y)\in\cX\times\cY$.
  \item A \emph{generative model} is a model for the joint probability
        $p(x,y)=\PP(X=x,Y=y)$ where $(x,y)\in\cX\times\cY$.
\end{enumerate}
\end{definition}

For instance, logistic regression is a discriminative model because we model
$\PP(Y=1\mid X=x)$ by $\phi(w\T x+b)$ and $\PP(Y=0\mid X=x)$ by
$1-\phi(w\T x+b)$. On the other hand, Naive Bayes classification is a
generative model because for $\cX=V^{p}$ we write $\PP(X=x,Y=y)$ as:
\begin{align}
  \PP(X=x,Y=y) &= \PP(X=x\mid Y=y)\,\PP(Y=y)\nonumber\\
               &\underset{\text{assumption}}{\overset{\text{Naive}}{\approx}}
                 \prod_{j=1}^{p}\PP(X^{j}=x_{j}\mid Y=y)\cdot\PP(Y=y).
  \label{eq:ch6-nb-joint}
\end{align}

Intuitively, a discriminative model directly models the difference between
data points with different labels, whereas a generative model first learns the
joint distribution $p(x,y)$ and then converts it into $p(y\mid x)$ for
classification. A discriminative model is often \emph{more efficient} than a
generative model. However, a generative model can be used to generate new data
similar to existing data, which is particularly useful in data-scarce
situations.

\subsection{Naive Assumption and Marginal Distribution Modeling}
\label{subsec:ch6-naive-marginals}

We explain first the discrete setting where the elementary objects are tuples
of size $p$:
\begin{equation}\label{eq:ch6-discrete-setting}
  \cX=\Omega^{p},\ \Omega\ \text{discrete}.
\end{equation}

The \emph{Naive} assumption is the following condition:
\begin{equation}\label{eq:ch6-naive-assumption}
  \forall (x_{1},\dots,x_{p})\in\cX,\ \forall y\in\cY:\
  \PP\bigl(\underbrace{(X^{1},\dots,X^{p})}_{=X}=(x_{1},\dots,x_{p})\,\big|\,Y=y\bigr)
  =\prod_{j=1}^{p}\PP(X^{j}=x_{j}\mid Y=y).
\end{equation}

\begin{warning}
This relation is to be understood from the perspective
of using a classifier, i.e., compare estimations for each label and choose
the largest one (this works well) and not as a prediction of the real
probabilities (does not work well); this is often formulated as:
\begin{center}
  \fbox{`naive Bayes is a good classifier but a bad predictor'}
\end{center}
\end{warning}

\begin{warning}
The assumption~\eqref{eq:ch6-discrete-setting} is taken to simplify the
setting. Note that in many cases it is satisfied only after some adjustments
of the dataset, for instance the dataset in Table~\ref{tab:ch6-training} does
not satisfy it because while most objects have 3 words, the second to last,
``amazing film'' has only 2 words. Several solutions can be proposed:
\begin{itemize}
  \item define $\cX=\bigcup_{p\ge 0}\Omega^{p}$ and assume that the length
        $P$ of an input object is a random variable that is
        \textbf{independent} of the label. Conditional to $P=p$ the
        relation~\eqref{eq:ch6-naive-assumption} is assumed true;
  \item or we can pad to the right with a special word called
        ``end-of-sequence'' token ``\texttt{<EOS>}'' to have
        $\Omega=V\cup\{\texttt{<EOS>}\}$, $p=3$ (or maximum value in general)
        and $\cX=\Omega^{p}$ (see definition of $V$
        in~\eqref{eq:ch6-vocab}).\footnote{In this case the model is not
        realistic because by definition no other word can occur after
        \texttt{<EOS>} so the naive Bayes independence assumption is not true
        in practice.}
\end{itemize}
In any case, we assume the model is valid even if this is questionable
from the modelisation point of view.
\end{warning}

Note that relation~\eqref{eq:ch6-discrete-setting} is not necessary a
`Bag-of-Words' model as in Section~\ref{sec:ch6-bow-example}. It becomes one
only after~\eqref{eq:ch6-naive-assumption} is assumed.

Indeed, under this assumption, given $Y$ (the label), we model the
(conditional) distribution of $X$ as the product of the (conditional) marginal
distributions of $X^{j}$, $1\le j\le p$. Such a product is independent of the
order of the factors.

But one may chose from the start some base event set that \textbf{intrinsically}
contains the Bag-of-Words assumption. That is, we may want to have elementary
events that only keep track of the number of times each `word' (possible
value) appears. So for $M=|\Omega|$ the elementary events are:
\begin{equation}\label{eq:ch6-bow-space}
  \mathcal{BOW}_{p}^{\Omega}:=\left\{(n_{1},\dots,n_{M})\in\N^{p}
    \,\middle|\,\sum_{m}n_{m}=p\right\}.
\end{equation}

Now, the question is: what distribution should be chosen on
$\mathcal{BOW}_{p}^{\Omega}$. We will give below a choice for the discrete
version and one for the continuous version.

\paragraph{$\bullet$ Categorical features ($\Omega$ discrete): Multinomial
distribution}

The multinomial distribution is a generalization of the binomial
distribution. Recall that a multinomial distribution models the following
phenomenon: we have an event with $M=|\Omega|$ possible values and denote:
\[
  \PP(\text{``$m$-th value''})=p_{m},\ 1\le m\le M.
\]
We draw $p$ times and let $W_{m}$, $1\le m\le M$ be the number of times the
$m$-th value appears. Then, for all
$\mathbf{n}=(n_{1},\dots,n_{M})\in\mathcal{BOW}_{p}^{\Omega}$, we have:
\[
  \PP(W_{1}=n_{1},\dots,W_{M}=n_{M})
   =\frac{p!}{n_{1}!\cdots n_{M}!}\,p_{1}^{n_{1}}\cdots p_{M}^{n_{M}}.
\]

For example in the previous section we constructed two multinomial
distributions: one for the Positive ($+$) sentiment and one for the Negative
($-$) sentiment. The set of events is $V$ defined in~\eqref{eq:ch6-vocab} with
$12$ values; here $p_{1}=\PP(\text{``love''})$, $p_{2}=\PP(\text{``this''})$
and so on.

Given the multinomial law on $\mathcal{BOW}_{p}^{\Omega}$, we may want now to
generate objects in $\cX$. The (conditional) generation of an object in the
set $\cX$ is compatible with the multinomial law on $\cX$; for instance, the
first item in the dataset, ``love this movie'' corresponds to
$n_{1}=n_{2}=n_{3}=1$ and $n_{k}=0$ for $k\ge 4$. Accordingly
\begin{align}
  \PP[X=\text{``love this movie''}\mid\text{Positive}]
    &= C\cdot(p_{1})^{1}(p_{2})^{1}(p_{3})^{1}(p_{4})^{0}\cdots(p_{12})^{0}
     \nonumber\\
    &\propto \PP[\text{``love''}\mid\text{Positive}]
            \cdot\PP[\text{``this''}\mid\text{Positive}] \nonumber\\
    &\qquad\cdot\PP[\text{``movie''}\mid\text{Positive}].
  \label{eq:ch6-multinom-example}
\end{align}

\paragraph{$\bullet$ Continuous features: Gaussian distribution}

In practice, a continuous feature is often modeled using a Gaussian (possibly
multivariate) distribution. For example, if $X=\texttt{patient height}$, we
can model it using a $\cN(\mu,\sigma^{2})$ distribution with estimates:
\begin{equation}\label{eq:ch6-gaussian-estimates}
  \mu\simeq \frac{1}{N}\sum_{i=1}^{N}X_{i}
  \quad\text{and}\quad
  \sigma\simeq \frac{1}{N}\sum_{i=1}^{N}(X_{i}-\mu)^{2}
\end{equation}
where $\{(X_{1},Y_{1}),\dots,(X_{N},Y_{N})\}$ is the training data. Note that
in this case, in the equation~\eqref{eq:ch6-discrete-setting} the density is
used instead (otherwise the result is always null).

\begin{remark}
More general situations can appear when each dimension has different
characteristics, for instance, some can be binary, others multinomial, and
others continuous. In this case, the state space is not $\Omega^{p}$ but of
type $\prod_{j=1}^{p}\Omega_{j}$ with $\Omega_{j}$ adapted to the dimension
$j$. The independence assumption helps in constructing the overall probability
as before.
\end{remark}

\begin{remark}
Naive Bayes has some similarities with logistic regression, for instance, the
presence of a linear decision boundary in the log-probability space. Both are
related to the log-odd ratio
$\log\!\left(\dfrac{\PP[Y=1\mid X]}{\PP[Y=0\mid X]}\right)$. But the naive
Bayes has a generative perspective as it is concerned with the joint
data-label distribution while logistic regression is a discriminative model
because it is concerned with the conditional probability. Naive Bayes can be
interpreted as a special case of Logistic Regression where the parameters
$w_{i}$ are determined by generative probabilities rather than optimization.
\end{remark}

\begin{tcolorbox}[advancedbox]
\textbf{Important technique 6.8.} \emph{Logarithms are commonly applied to
probability computations to improve numerical stability. Given that the
posterior probability is expressed as}
\begin{equation}\label{eq:ch6-log-posterior-product}
  \PP(Y\mid X)\propto \PP(Y)\prod_{i}\PP(X_{i}\mid Y),
\end{equation}
\emph{taking the logarithm transforms the product of probabilities into a
summation:}
\begin{equation}\label{eq:ch6-log-posterior-sum}
  \log\PP(Y\mid X)\propto \log\PP(Y)+\sum_{i}\log\PP(X_{i}\mid Y).
\end{equation}
\emph{This conversion mitigates the effects of numerical underflow when dealing
with very small probability values.}
\end{tcolorbox}

\section{Exercises: Naive Bayes}
\label{sec:ch6-exercises}

\begin{tcolorbox}[remarkbox]
\textbf{Exercise 6.1} (mixed formulation)\textbf{.} Formulate naive Bayes when
$p=2$ and the first feature (dimension) is discrete and the second feature
continuous: use the product between probability and density and explain the
relationship with likelihood.
\end{tcolorbox}

\section{Implementations: naive Bayes}
\label{sec:ch6-implementations}

\begin{tcolorbox}[remarkbox]
\textbf{Exercise 6.2} (Naive Bayes on MNIST)\textbf{.}
\begin{enumerate}
  \item Consider the MNIST dataset. Binarize it, that is, transform each
        image into an element of $\{0,1\}^{784}$. Compute accuracy of the
        naive Bayes classification.
  \item Do the same on CIFAR-10 after a transformation to grayscale and
        binarization.
\end{enumerate}
Note: for CIFAR-10 an alternative form is to use quantization instead of
binarization.
\end{tcolorbox}
